% Applied Software Development Lab -- Lab Handout 7: At the end, we shoot and clean up
% Copyright (c) Michael Lehn. Licensed under CC BY-SA 4.0.
%
% Region coordinates (bp from the top of the page) follow the page layout of
% the original German handout "Laborblatt 7".
\documentclass{aslab}
\usepackage{unicode-math}
\setmathfont{texgyretermes-math.otf}[Scale=1.2]
\newfontfamily\figserif{TeX Gyre Termes}
\sethandout{7}{\texorpdfstring{At the end, we shoot \emoji{1F603}\
  {\fontsize{8bp}{9bp}\selectfont and clean up \emoji{1F629}}}%
  {At the end, we shoot and clean up}}
\graphicspath{{../figures/handout07/}}

%% Python listings with the color scheme of the original
%% inline code in bold lines (the original keeps it in regular weight)
\newcommand*\codemd[1]{{\mdseries\codearg{#1}}}
\newcommand*\cy[1]{\textcolor{codecyan}{#1}}
\makeatletter
\newcommand*\regionpitch{\baselineskip=\aslab@pitch\relax}
\makeatother
\lstdefinestyle{py}{style=aslab,sensitive=true,escapechar=|,
  basicstyle=\ttfamily\regionpitch,
  morekeywords=[1]{class,def,if,elif,else,for,while,or,and,in,return,not},
  keywordstyle=[1]\color{codeorange},
  morekeywords=[2]{import,from},keywordstyle=[2]\color{codemagenta},
  morekeywords=[3]{None,True,False,range,len},keywordstyle=[3]\color{codecyan},
  morecomment=[l]\#,commentstyle=\color{codepurple},
  alsoletter={0123456789},emphstyle=\color{codered},
  emph={0,1,2,3,4,5,6,7,8,9,01,10,20,30,80,100,105,165,200,223,255,300,400,
    570,600,800}}
%% numbered listing (11bp): "nnn code" with the numbers in orange
\lstdefinestyle{pyn}{style=py,numbers=left,
  numberstyle=\ttfamily\color{codeorange},numbersep=6.6bp,xleftmargin=26.4bp}
%% numbered listing (10bp)
\lstdefinestyle{pynsmall}{style=py,numbers=left,
  basicstyle=\fontsize{10bp}{11bp}\selectfont\ttfamily\regionpitch,
  numberstyle=\fontsize{10bp}{11bp}\selectfont\ttfamily\color{codeorange},
  numbersep=6bp,xleftmargin=24bp}
\lstdefinestyle{pysmall}{style=py,
  basicstyle=\fontsize{10bp}{11bp}\selectfont\ttfamily}

%% lists of this handout
\newlist{ntasks}{enumerate}{1}
\setlist[ntasks]{label={\arabic*.},leftmargin=18bp,labelsep=*,align=left,
  labelindent=0pt}
\newlist{alist}{enumerate}{1}
\setlist[alist]{label={(\Alph*)},leftmargin=18bp,labelsep=*,align=left,
  labelindent=0pt}
\newlist{dotlist}{itemize}{1}
\setlist[dotlist]{label={\textbullet},leftmargin=9bp,labelsep=*,align=left,
  labelindent=0pt}
\newcommand*\bl[1]{\textcolor{aslabblue}{#1}}
\lineskiplimit=-\maxdimen
\normallineskiplimit=-\maxdimen

\begin{document}
%% ===================================================================== 1 ==
\maketitleblock

\begin{region}{130.5}
\heading{Learning objectives}
\end{region}

\begin{region}[pitch=13]{145.0}
{\color{aslabblue}\bfseries
\begin{dashlist}
\item Deepening object-oriented programming
\item Refactoring existing programs
\item Incremental software development: small changes, testing, and versioning
\end{dashlist}\par}
\end{region}

\begin{region}{198.0}
\heading{Step 1: Cloning your own repository}
\end{region}

\begin{region}[pitch=13.87]{225.6}
Today we are working ``only'' with Python again. The goal is to further
improve the existing code (keyword: \emph{refactoring}) and then to extend its
functionality.\\
Before we start making changes, please take another look at the current
state of the game.
\end{region}

\begin{region}{295.7}
\textbf{As always:} Open two terminals -- one for editing (left), one for
running commands (right).
\end{region}

\begin{region}[pitch=13.3]{322.0}
\subheading{Left terminal (editing)}
Clone the repository and change into the working directory:
\end{region}

\begin{region}[pitch=13.2]{363.0}
\code{git clone https://gitlab.uni-ulm.de/USERNAME/asp-ulm}\\
\code{cd asp-ulm/python}
\end{region}

\begin{region}{403.3}
\subheading{Task 1:}
\end{region}

\begin{region}[pitch=13.3]{431.1}
Open the file \code{breakout.py} in your editor and examine the structure of
the program.\\
Focus on the following points:
\end{region}

\begin{region}[pitch=13.35]{457.0}
\begin{dashlist}
\item As a group, go through the flow of the program together, in particular
  the main loop at the end of the file.
\item Get an overview of the existing classes and their responsibilities.
\item Look specifically for global variables (e.g., \code{screen},
  \code{width}, \code{height}, \code{balls}, \code{running}) and determine
  where they are used or modified.
\item Distinguish which classes are instantiated (\code{Ball}) and which are
  only used as a class object with static attributes and methods
  (\code{Paddle}).
\item In doing so, recap the role of \code{__init__} and \code{self} in
  classes with instances.
\end{dashlist}
\end{region}

\begin{region}[pitch=13.2]{577.4}
Talk to each other until everyone has understood the basic structure. This
preparation is important for the refactoring in the next step.
\end{region}

\begin{region}[pitch=13.3]{619.0}
\subheading{Right terminal (running commands)}
Change into the Python directory here as well:
\end{region}

\begin{region}{658.9}
\code{cd asp-ulm/python}
\end{region}

\begin{region}{686.0}
\subheading{Task 2}
\end{region}

\begin{region}{712.7}
Run the game with the command \code{python breakout.py}.
\end{region}

\begin{region}[pitch=13.65]{739.1}
Compare the observed behavior with what you found out about the code in
Task~1. If there are any discrepancies or open questions, discuss them with us
before you continue.
\end{region}

\nextpage
%% ===================================================================== 2 ==
\begin{region}{69.4}
\heading{Step 2: Refactoring}
\end{region}

\begin{region}[pitch=13.3]{97.0}
The previous global variables are moved into a separate class Game. This
class goes into a new file \code{game.py}. From now on, all other classes are
also each moved into a file of their own, so that each file contains exactly
one piece of functionality and the code remains easy to follow.
\end{region}

\begin{region}{164.2}
\subheading{The goal is:}
\end{region}

\begin{region}[pitch=13.8]{177.0}
\begin{dashlist}
\item better structure
\item fewer global variables \textbf{outside of classes}
\item preparation for further object-oriented steps
\end{dashlist}
In the end, the file \code{breakout.py} should contain little more than the
main loop.
\end{region}

\begin{region}[pitch=13.6]{245.7}
\subheading{Task 1}
Create the file \code{game.py} with the following content:
\end{region}

\begin{region}[pitch=11]{285.2}
\begin{lstlisting}[style=pynsmall]
import pygame

class |\cy{Game}|:
    width, height = None, None
    screen = None

    def |\cy{init}|():
        Game.width, Game.height = 800, 600
        pygame.init()
        Game.screen = pygame.display.set_mode((Game.width, Game.height))

    def |\cy{quit}|():
        pygame.quit()
\end{lstlisting}
\end{region}

\begin{region}[pitch=13.6]{456.0}
\textbf{Task 2: \codemd{ball.py} (reference for the file at the end of the
task)}\\
Create a new file \code{ball.py} by copying: \code{cp breakout.py ball.py}
\end{region}

\begin{region}{497.4}
Edit the new file \code{ball.py} as follows:
\end{region}

\begin{region}[pitch=13.55]{510.3}
\begin{dashlist}
\item Delete everything except the class \code{Ball}.\\
  Tip for \textbf{nano}: Hold down the \textbf{Shift key} and use the arrow
  keys to select text.\\
  Press \textbf{Control+K} to delete the selected text.
\item Add the necessary \code{import} lines at the beginning:
\begin{lstlisting}[style=pyn]
from game import Game
import pygame
\end{lstlisting}
  Note:\\
  With \code{from game import Game}, the class Game is imported directly from
  the module \code{game.py}. This means you can simply write
  \code{Game.width} or \code{Game.screen} in the code.\\
  If you used \code{import game} instead, you would have to include the module
  name every time, e.g., \code{game.Game.width}.\\
  Both variants are correct -- but here we choose
  \code{from ... import ...} to keep the code easier to read.
\item Adjust how global variables are accessed: instead of \code{width},
  \code{height}, or \code{screen}, you now use \code{Game.width},
  \code{Game.height}, \code{Game.screen}.
\end{dashlist}
\end{region}

\nextpage
%% ===================================================================== 3 ==
\begin{region}{69.4}
\textbf{Reference: \codemd{ball.py}}
\end{region}

\begin{region}{97.3}
\begin{lstlisting}[style=pyn]
from game import Game
import pygame

class |\cy{Ball}|:
    def |\cy{\_\_init\_\_}|(self, x, y, vx, vy):
        self.x = x
        self.y = y
        self.vx = vx
        self.vy = vy
        self.c = (255, 255, 0)

    def |\cy{move}|(self):
        self.x += self.vx
        self.y += self.vy
        if self.x < 0 or self.x > Game.width:
            self.vx *= -1
        if self.y < 0 or self.y > Game.height:
            self.vy *= -1

    def |\cy{draw}|(self, surface):
        pygame.draw.circle(surface, self.c, (self.x, self.y), 5)
\end{lstlisting}
\end{region}

\begin{region}[pitch=13.3]{363.2}
\textbf{Task 2: \codemd{paddle.py}}\\
Move the class Paddle into a separate file \code{paddle.py} -- analogous to
\code{ball.py}.\\
Make sure that here, too, global variables are accessed via the class Game.\\
The position is now no longer initialized directly in the definition, but by
calling \code{Paddle.init()}.
\end{region}

\begin{region}{444.0}
\textbf{Reference: \codemd{paddle.py}}
\end{region}

\begin{region}{470.7}
\begin{lstlisting}[style=pyn]
from game import Game
import pygame

class |\cy{Paddle}|:
    c = (0, 0, 255)
    w, h = 100, 30
    x, y = None, None

    def |\cy{init}|():
        Paddle.x, Paddle.y = (Game.width - Paddle.w) // 2, 570

    def |\cy{left}|():
        Paddle.x -= 10
        if Paddle.x < 0:
            Paddle.x = 0

    def |\cy{right}|():
        Paddle.x += 10
        if Paddle.x + Paddle.w > Game.width:
            Paddle.x = Game.width - Paddle.w

    def |\cy{draw}|(surface):
        pygame.draw.rect(surface, Paddle.c,
                         (Paddle.x, Paddle.y, Paddle.w, Paddle.h))
\end{lstlisting}
\end{region}

\nextpage
%% ===================================================================== 4 ==
\begin{region}[pitch=13.45]{69.4}
\subheading{Task 3}
Clean up the file \code{breakout.py} so that it contains only the game flow
in the main loop. Import the required classes from the files \code{game.py},
\code{ball.py}, and \code{paddle.py}. Again, adjust the access to global
variables. At the beginning, initialize the game with \code{Game.init()} and
then call \code{Paddle.init()}.
\end{region}

\begin{region}{150.6}
\textbf{Reference: \codemd{breakout.py}}
\end{region}

\begin{region}{164.2}
\begin{lstlisting}[style=pyn]
from game import Game
from ball import Ball
from paddle import Paddle
import time
import pygame

Game.init()
Paddle.init()

running = True
balls = [Ball(400, 300, -1, 2), Ball(200, 10, 4, 2)]
while running:
    for event in pygame.event.get():
        if event.type == pygame.QUIT:
            running = False
        elif event.type == pygame.KEYDOWN:
            if event.key == pygame.K_q:
                running = False
    keys = pygame.key.get_pressed()
    if keys[pygame.K_a]:
        balls.append(Ball(400, 300, -1, 2))
    if keys[pygame.K_LEFT]:
       Paddle.left()
    if keys[pygame.K_RIGHT]:
       Paddle.right()

    pygame.Surface.fill(Game.screen, (0, 0, 0))
    for i in range(0,len(balls)):
        balls[i].draw(Game.screen)
        balls[i].move()
    Paddle.draw(Game.screen)
    pygame.display.flip()
    time.sleep(1/30)

Game.quit()
\end{lstlisting}
\end{region}

\begin{region}[pitch=13.37]{598.0}
\subheading{Task 4: Functional test, git commit, and push}
Start the game as usual with \code{python breakout.py}. Check whether the game
works as before. To do so, move the paddle all the way to the left and all the
way to the right at least once. This runs through all relevant conditions in
\code{Paddle.left()} and \code{Paddle.right()}. If, for example, you forgot to
replace \code{width} with \code{Game.width} during refactoring, you will
notice it now at the latest. Of course, this test is not complete -- but it
covers typical errors that can occur during restructuring.
\end{region}

\begin{region}[pitch=13.2]{704.5}
If the game works as expected, add the new files to the repository and make a
commit:\\
\code{git add ball.py paddle.py game.py}\\
\code{git commit -a -m "one file does one thing"}\\
\code{git push}
\end{region}

\nextpage
%% ===================================================================== 5 ==
\begin{region}{69.4}
\heading{Step 3: Simplifying interfaces}
\end{region}

\begin{region}{97.3}
The method \code{draw} of the class \code{Ball} (in \code{ball.py}) currently
looks like this:
\end{region}

\begin{region}{124.0}
\begin{lstlisting}[style=py]
# ...
\end{lstlisting}
\end{region}

\begin{region}{137.2}
\begin{lstlisting}[style=py]
class |\cy{Ball}|:
    # ...

    def |\cy{draw}|(self, surface):
        pygame.draw.circle(surface, self.c, (self.x, self.y), 5)
\end{lstlisting}
\end{region}

\begin{region}{210.3}
It is called in only one place -- in the main loop of \code{breakout.py}:
\end{region}

\begin{region}{237.0}
\begin{lstlisting}[style=py]
    balls[i].draw(Game.screen)
\end{lstlisting}
\end{region}

\begin{region}[pitch=13.5]{264.4}
So drawing always happens on \code{Game.screen}. That makes sense -- we have
exactly one \textbf{game} and one \textbf{drawing surface}.\\
Passing the drawing surface is therefore redundant and makes the code
needlessly general.\\
Our code will become complex soon enough. So: out with everything that is not
needed yet.
\end{region}

\begin{region}{345.6}
\subheading{Task:}
\end{region}

\begin{region}[pitch=13.3]{372.4}
\begin{ntasks}
\item Simplify the method \code{draw} in the class \code{Ball} as follows:
\begin{lstlisting}[style=py]
# ...
class |\cy{Ball}|:
  # ...

   def |\cy{draw}|(self):
        pygame.draw.circle(Game.screen, self.c, (self.x, self.y), 5)
\end{lstlisting}
  \blank
  Make sure that you now use Game.screen directly in the method draw.\\
  Also adjust the call in \code{breakout.py} accordingly:
  \code{balls[i].draw()}
\end{ntasks}
\end{region}

\begin{region}[pitch=13.53]{505.6}
\begin{ntasks}[start=2]
\item \textbf{Transfer task:}\\
  The method \code{draw} of the class Paddle has the same unnecessary
  flexibility.\\
  Adjust it in the same way so that it uses \code{Game.screen} directly --
  and can then simply be called in \code{breakout.py} with
  \code{Paddle.draw()}.
\end{ntasks}
\end{region}

\begin{region}[pitch=13.45]{559.5}
\begin{ntasks}[start=3]
\item If everything works as expected, commit your changes:\\
  \code{git commit -a -m "simplified draw methods"}\\
  \code{git push}
\end{ntasks}
\end{region}

\nextpage
%% ===================================================================== 6 ==
\begin{region}{69.4}
\heading{Step 4: Reflecting the ball off the walls}
\end{region}

\begin{region}[pitch=14.2]{97.0}
It is high time to show a sketch of our playing field with coordinates:\\
As you may have already noticed (or are noticing now), the origin is at the
\textbf{top left}, and the \textbf{y-axis points downward} -- as is customary
in Pygame.
\end{region}

\begin{region}[pitch=13.2]{151.9}
The following sketch shows two balls, each moving toward a wall. Without
reflection, they would pass through the wall in the next frame. The reflection
prevents this -- by reversing the direction of motion along the respective
axis.
\end{region}

\begin{region}{205.5}
\subheading{The sketch shows:}
\end{region}

\begin{region}[pitch=13.47]{218.5}
\begin{dashlist}
\item the current center of the ball at \code{(self.x, self.y)}
\item the new, unreflected position \code{(new_x, new_y)}
\item the corresponding velocity vector \code{(vx, vy)}
\item the impact on the left or right wall and the corresponding reflection
\end{dashlist}
\end{region}

\placefigure{68.7}{263.2}{481.9}{wall-reflection.tex}

\begin{region}{490.8}
Adapt the \code{move} method in the class \code{Ball} (in \code{ball.py}) as
follows:
\end{region}

\begin{region}{517.4}
\begin{lstlisting}[style=pyn,firstnumber=13]
    def |\cy{move}|(self):
\end{lstlisting}
\end{region}

\begin{region}{530.6}
\begin{lstlisting}[style=pyn,firstnumber=14]
        new_x = self.x + self.vx
        new_y = self.y + self.vy

        if new_x < 0:
            new_x = -new_x
            self.vx = -self.vx
        if new_x > Game.width:
            new_x = 2 * Game.width - new_x
            self.vx = -self.vx
        if new_y < 0:

            self.vy = -self.vy
        if new_y > Game.height:

            self.vy = -self.vy

        self.x = new_x
        self.y = new_y
\end{lstlisting}
\end{region}

\nextpage
%% ===================================================================== 7 ==
\begin{region}{83.7}
\subheading{Tasks}
\end{region}

\begin{region}[pitch=13.2]{111.6}
\begin{ntasks}
\item Start the game with \code{python breakout.py} and observe how the
  reflection of the ball at the left and right walls has changed.
\end{ntasks}
\end{region}

\begin{region}[x=74.7,right=0,pitch=14.15]{151.2}
You should see that the ball (more precisely: its center) is reflected
realistically. Even at high speed, it always stays \textbf{inside the drawing
surface} -- provided we regard the ball as a \textbf{point mass}. The ball
itself can therefore stick halfway into the wall.
\end{region}

\begin{region}[x=74.7,right=0]{207.0}
\textbf{How do you arrive at the code for the reflection?}
\end{region}

\begin{region}[x=74.7,right=0,pitch=17.1]{223.6}
\begin{alist}
\item If \code{new_x < 0}, the ball has penetrated the left wall by
  \bl{$-$new\_x}. With a reflection, the new x-coordinate lies symmetrically
  to \bl{$0$} at \bl{$-$new\_x}.
\end{alist}
\end{region}

\begin{region}[x=74.7,right=0,pitch=14.6]{257.8}
\begin{alist}[start=2]
\item If \code{new_x > Game.width}, the ball has moved beyond the right edge
  in the x-direction by \bl{new\_x $-$ Game.width}. The reflected
  x-coordinate lies symmetrically to \bl{Game.width} at\\
  \bl{Game.width $-$ (new\_x $-$ Game.width) $=$ $2\cdot{}$Game.width $-$
  new\_x}
\end{alist}
\end{region}

\begin{region}[x=74.7,right=0,pitch=13.3]{316.1}
\textbf{Note on how to improve the reflection (optional):}\\
If you want to be really precise, the ball would have to bounce off at
\end{region}

\begin{region}[x=74.7,right=0,pitch=14.3]{343.5}
\begin{ntasks}[start=3]
\item \code{x = 5} \textbf{(left wall)}
\item \code{x = Game.width - 5} \textbf{(right wall)}
\end{ntasks}
\end{region}

\begin{region}[x=74.7,right=0]{371.1}
(assuming that the radius of the ball is 5).
\end{region}

\begin{region}[x=74.7,right=0,pitch=14.2]{398.4}
You should also observe that the ball can still \textbf{leave the screen at
the top and bottom} -- we will fix that next.
\end{region}

\begin{region}[pitch=13.8]{440.1}
\begin{ntasks}[start=2]
\item \textbf{Transfer task: Reflection at the floor and ceiling}\\
  Adjust the code in \textbf{line 24} and \textbf{line 27} so that the ball is
  also reflected realistically at the top and bottom -- analogous to the side
  walls.
\end{ntasks}
\end{region}

\begin{region}[pitch=13.37]{494.0}
\begin{ntasks}[start=3]
\item If the game works correctly and the ball now also reflects off the
  ceiling and floor, commit your changes:\\
  \code{git commit -a -m "better reflection"}\\
  \code{git push}
\end{ntasks}
\end{region}

\nextpage
%% ===================================================================== 8 ==
\begin{region}{69.4}
\heading{Step 5: Reflecting the ball off the paddle}
\end{region}

\begin{region}[pitch=13.67]{97.0}
The following sketch shows two situations in which a ball can hit the
paddle.\\
For each ball, the sketch shows its current position \code{(self.x, self.y)}
and the new position \code{(new_x,new_y)} calculated for the case that the ball \textbf{keeps flying
unhindered} -- that is, \textbf{without collision}, as is still the case in
the current code:
\end{region}

\begin{region}[pitch=14.5]{151.2}
\begin{dashlist}
\item The ball on the \textbf{left} is initially located to the left of the
  paddle. Its motion vector crosses the top edge of the paddle, and the new
  position lies \textbf{below the edge and inside the paddle}.
\item The ball on the \textbf{right} is initially located \textbf{above} the
  paddle. Its new position lies \textbf{to the right of the paddle}, but also
  \textbf{below the top edge}.
\end{dashlist}
\end{region}

\placefigure{56.7}{213.2}{450.8}{paddle-impact.tex}

\begin{region}[pitch=14.2]{346.7}
In both cases: \textbf{the motion vector intersects the top edge of the
paddle}. The reflection should take place at this \textbf{impact point}.
\end{region}

\begin{region}[pitch=13.2]{387.4}
To calculate this intersection point as simply and generally as possible, let
us look at the following sketch:
\end{region}

\placefigure{56.7}{417.6}{266.5}{similar-triangles.tex}

\begin{region}[pitch=13.6]{618.5}
In the sketch you can see two similar triangles -- an ideal case for the
\textbf{good old intercept theorem}. It is simple, intuitive, and leads
directly to the solution -- and yet it is often forgotten at university or
replaced by needlessly complicated calculations. This shows how helpful solid
school geometry can be in programming, too.
\end{region}

\begin{region}[pitch=13.25]{686.6}
Our goal: we want to calculate the \textbf{x-coordinate of the impact point}
-- that is, the point at which the motion vector of the ball intersects the
top edge of the paddle. To do this, we set up the ratios of corresponding
sides of the two triangles:
\end{region}

\begin{region}[x=147.6,width=390]{747.5}
$\dfrac{\text{Dx}}{\text{Dy}} = \dfrac{\text{dx}}{\text{dy}}$\ \ with\ \
$\text{dx} = \text{new\_x - self.x}$\ \ and\ \
$\text{dy} = \text{new\_y - self.y}$
\end{region}

\nextpage
%% ===================================================================== 9 ==
\begin{region}{81.7}
The vertical distance to the paddle edge is:
\end{region}

\begin{region}{111.3}
\centering $\text{Dy} = \text{Paddle.y} - \text{self.y}$
\end{region}

\begin{region}{138.3}
Then the horizontal component can be calculated with:
\end{region}

\begin{region}{172.7}
\centering $\text{Dx} = \dfrac{\text{dx}}{\text{dy}}\text{Dy}$
\end{region}

\begin{region}{197.1}
Finally, the x-coordinate of the impact point is:
\end{region}

\begin{region}{226.7}
\centering $\text{px} = \text{self.x} + \text{Dx}$
\end{region}

\begin{region}{254.6}
\subheading{Tasks}
\end{region}

\begin{region}[pitch=13.75]{282.1}
\begin{ntasks}
\item The following code fragment checks whether the ball crosses the
  \textbf{line through the top edge of the paddle}. If so, the
  \textbf{intersection point} with this edge is to be calculated (you fill in
  this gap in the code yourself). Then it checks whether the intersection
  point lies within the \textbf{horizontal extent of the paddle}. If so, the
  ball is reflected by reversing its vertical velocity.
\end{ntasks}
\end{region}

\begin{region}[x=74.7,right=0,pitch=13.2]{363.7}
\begin{lstlisting}[style=py]
if self.y < Paddle.y and new_y >= Paddle.y:
    # Calculate px in this gap
    if px >= Paddle.x and px <= Paddle.x + Paddle.w:
        new_y = Paddle.y + Dy - dy
        self.vy = -self.vy
\end{lstlisting}
\end{region}

\begin{region}[x=74.7,right=0,pitch=13.3]{443.9}
\textbf{Transfer task:} Think about where in the code this fragment should be
inserted so that the ball is reflected when it hits the paddle.
\end{region}

\begin{region}[x=74.7,right=0,pitch=13.25]{484.5}
\textbf{Note for the style-conscious:} Use \textbf{integer division}
(\code{//}) wherever it makes sense. Even if it does not matter in this small
game -- deep down, every HPC programmer's heart aches a little at unnecessary
floating-point division.
\end{region}

\begin{region}[pitch=13.45]{537.3}
\begin{ntasks}[start=2]
\item If the ball now reflects off the paddle as intended, commit your
  changes:\\
  \code{git commit -a -m "paddle does what it should"}\\
  \code{git push}
\end{ntasks}
\end{region}

\nextpage
%% ==================================================================== 10 ==
\begin{region}{69.4}
\heading{Step 6: Breakout-style bouncing off the paddle}
\end{region}

\begin{region}[pitch=13.8]{96.9}
In classic Breakout games, the ball is \textbf{not simply reflected} when it
hits the paddle. Instead, its new direction of motion depends on \textbf{where
on the paddle} the ball hits.
\end{region}

\begin{region}{150.8}
A simple model:
\end{region}

\begin{region}[pitch=14.3]{177.1}
\begin{dashlist}
\item If the ball hits the paddle \textbf{exactly in the middle}, it flies
  \textbf{straight up}.
\item If it hits the \textbf{left edge}, it flies \textbf{up and to the
  left}.
\item If it hits the \textbf{right edge}, it flies \textbf{up and to the
  right}.
\end{dashlist}
\end{region}

\begin{region}[pitch=14.2]{234.1}
The ball should keep its previous \textbf{speed (magnitude)}. Only the
\textbf{direction} changes -- and it does so \textbf{linearly depending on the
impact point}.
\end{region}

\begin{region}{274.8}
The following sketch illustrates this idea:
\end{region}

\placefigure{56.7}{278.7}{429.8}{breakout-bounce.tex}

\begin{region}{478.8}
We describe the \textbf{impact point} on the paddle by a parameter
$t\in[-1,1]$ such that:
\end{region}

\begin{region}[pitch=17.15]{495.9}
\begin{dashlist}
\item $t=-1$: ``far left'',
\item $t=0$: ``exactly in the middle'',
\item $t=1$: ``far right''.
\end{dashlist}
\end{region}

\begin{region}{544.2}
To do so, we parametrize the x-coordinate of the impact point \code{px} by
\end{region}

\begin{region}{565.5}
\centering $\text{px} = \text{Paddle.x} + \dfrac{\text{Paddle.w}}{2}\cdot(1+t)
  \quad\text{with}\quad t\in[-1,1]$
\end{region}

\begin{region}{587.7}
Solving for \code{t} yields:
\end{region}

\begin{region}{610.1}
\centering $t = (\text{px} - \text{Paddle.x})\cdot\dfrac{2}{\text{Paddle.w}}-1$
\end{region}

\begin{region}[pitch=13.8]{635.4}
Now we can map the parameter $t\in[-1,1]$ to an angle $\varphi$ in
\textbf{radians}. A few annoying details have to be taken into account -- for
example, that in our coordinate system the \textbf{y-axis points downward}.
The following formula already takes this into account correctly:
\end{region}

\begin{region}{700.5}
\centering $\varphi = \left(\dfrac{t}{3}-1\right)\cdot\dfrac{\pi}{2}\ \in
  \left[-\dfrac{4\pi}{6}, -\dfrac{2\pi}{6}\right].$
\end{region}

\begin{region}{725.7}
This means:
\end{region}

\begin{region}[pitch=17.15]{742.2}
\begin{dashlist}
\item $t=-1$ gives an angle of \textbf{60° up to the left.}
\item $t=1$ gives \textbf{60° up to the right.}
\item $t=0$ gives \textbf{straight up.}
\end{dashlist}
\end{region}

\nextpage
%% ==================================================================== 11 ==
\begin{region}{81.7}
The magnitude of the velocity is preserved:
\end{region}

\begin{region}{103.3}
\centering $v = \sqrt{\text{self.vx}^2 + \text{self.vy}^2}$
\end{region}

\begin{region}{122.8}
The new velocity vector is then given by:
\end{region}

\begin{region}{139.3}
\centering $\text{self.vx} = v\cdot\cos\varphi$\quad and\quad
  $\text{self.vy} = v\cdot\sin\varphi$.
\end{region}

\begin{region}{180.4}
\subheading{Tasks}
\end{region}

\begin{region}{207.9}
\begin{ntasks}
\item \textbf{Now let's put the theory into practice again!}
\end{ntasks}
\end{region}

\begin{region}[x=74.7,right=0,pitch=13.5]{234.6}
For the implementation, we need some mathematical functions provided by the
module \code{math}. Import the module in \code{ball.py}:
\end{region}

\begin{region}[x=74.7,right=0]{274.8}
\begin{lstlisting}[style=py]
import math
\end{lstlisting}
\end{region}

\begin{region}[x=74.7,right=0]{301.2}
What we need from it:
\end{region}

\begin{region}[x=74.7,right=0]{318.4}
\begin{dotlist}
\item \code{math.pi}: the constant $\pi$
\end{dotlist}
\end{region}

\begin{region}[x=74.7,right=0]{343.2}
\begin{dotlist}
\item \code{math.hypot(x, y)}: computes $\sqrt{x^2+y^2}$
\end{dotlist}
\end{region}

\begin{region}[x=74.7,right=0]{361.0}
\begin{dotlist}
\item \code{math.cos} and \code{math.sin} compute the cosine and sine of an
  angle in \textbf{radians}.
\end{dotlist}
\end{region}

\begin{region}[x=74.7,right=0,pitch=14.1]{386.9}
With these tools, it should not be difficult to complete the following code
fragment and to place it \textbf{at a suitable point} in the code:
\end{region}

\begin{region}[x=74.7,right=0,pitch=13.2]{427.8}
\begin{lstlisting}[style=py]
if px >= Paddle.x and px <= Paddle.x + Paddle.w:
    new_y = Paddle.y + Dy - dy
    # Calculate t, phi, v, self.vx, self.vy in this gap
\end{lstlisting}
\end{region}

\begin{region}[x=74.7,right=0,pitch=13.75]{481.3}
\textbf{Remark:} Strictly speaking, we first still reflect the ball according
to the principle ``angle of incidence = angle of reflection'' -- and only then
change the direction according to the \textbf{Breakout rules}.\\
This is not entirely clean, but to the human eye it will look exactly as if
we reflected the ball Breakout-style right away.
\end{region}

\begin{region}{564.4}
\begin{ntasks}[start=2]
\item \textbf{If everything works, then as always: into the repository!}
\end{ntasks}
\end{region}

\begin{region}[x=74.7,right=0,pitch=13.5]{591.2}
\code{git commit -a -m "implemented Breakout reflection"}\\
\code{git push}
\end{region}

\nextpage
%% ==================================================================== 12 ==
\begin{region}{69.4}
\heading{Step 7: Something without theory for a change}
\end{region}

\begin{region}[pitch=13.5]{97.0}
Now we simply make the code a little nicer -- without any theory at all.\\
In \code{breakout.py}, replace this loop:
\end{region}

\begin{region}[pitch=12.6]{137.1}
\begin{lstlisting}[style=py]
    for i in range(0,len(balls)):
        balls[i].draw()
        balls[i].move()
\end{lstlisting}
\end{region}

\begin{region}{187.3}
with the more compact and more readable variant:
\end{region}

\begin{region}[pitch=12.6]{213.8}
\begin{lstlisting}[style=py]
    for b in balls:
        b.draw()
        b.move()
\end{lstlisting}
\end{region}

\begin{region}[pitch=13.1]{264.0}
Both variants do exactly the same thing -- but the second one is clearer,
shorter, and more in line with Python style.
\end{region}

\begin{region}[pitch=13.6]{304.4}
\textbf{As always:} Test whether everything still works as expected. Then
off to the repository with\\
\code{git commit -a -m "nicer loop"}\\
\code{git push}
\end{region}

\begin{region}{371.4}
\heading{Step 8: Quickly adding a new feature (rockets!)}
\end{region}

\begin{region}[pitch=13.2]{397.9}
We have now spent a lot of time cleaning up our code, structuring it more
nicely, and solving details such as the ball reflection properly. High time
to check whether all this pays off in practice!
\end{region}

\begin{region}[pitch=13.25]{451.5}
Our paddle now gets a new ability: \textbf{firing rockets!}\\
The rockets should fly upward, i.e., similar to balls, only in the opposite
direction.\\
At first, they are even displayed exactly like balls -- we may change that
later.
\end{region}

\begin{region}[pitch=13.5]{504.6}
Create a new file \code{rocket.py}. The easiest way is probably to copy the
file \code{ball.py} and adapt it: \code{cp ball.py rocket.py}\ \ Then replace
its content with:
\end{region}

\begin{region}{544.8}
\begin{lstlisting}[style=py]
from game import Game
\end{lstlisting}
\end{region}

\begin{region}{558.0}
\begin{lstlisting}[style=py]
from paddle import Paddle
import pygame
import math

class |\cy{Rocket}|:
    def |\cy{\_\_init\_\_}|(self, x, y, vy):
        self.x = x
        self.y = y
        self.vy = vy
        self.c = (255, 255, 255)

    def |\cy{move}|(self):
        if self.y < 10:
            return
        self.y += self.vy

    def |\cy{draw}|(self):
        pygame.draw.circle(Game.screen, self.c, (self.x, self.y), 5)
\end{lstlisting}
\end{region}

\nextpage
%% ==================================================================== 13 ==
\begin{region}{83.4}
\subheading{Tasks}
\end{region}

\begin{region}{107.8}
The following changes in \code{breakout.py} are sufficient:
\end{region}

\begin{region}{132.8}
\begin{ntasks}
\item \textbf{Importing the rocket class}
\end{ntasks}
\end{region}

\begin{region}[x=74.7,right=0]{157.1}
At the very top, with the imports (preferably next to the other
\code{from ... import ...} lines), add:
\end{region}

\begin{region}[x=74.7,right=0]{194.6}
\begin{lstlisting}[style=py]
from rocket import Rocket
\end{lstlisting}
\end{region}

\begin{region}[x=74.7,right=0]{219.7}
This lets you write \code{Rocket(...)} directly later on, instead of
\code{rocket.Rocket(...)}.
\end{region}

\begin{region}{244.7}
\begin{ntasks}[start=2]
\item \textbf{Creating the list of rockets}\\
  Near the place where the balls are initialized (e.g., directly after it),
  add:
\end{ntasks}
\end{region}

\begin{region}[x=74.7,right=0]{281.0}
\tcode{rockets = []}
\end{region}

\begin{region}[x=74.7,right=0]{305.0}
At first, everything is peaceful -- no rockets on the way yet.
\end{region}

\begin{region}[pitch=12.3]{330.4}
\begin{ntasks}[start=3]
\item \textbf{Catching the arrow key ↑ and creating a rocket}\\
  Extend the handling of key-press events (where you already check for
  \code{pygame.KEYDOWN} and \code{event.key == pygame.K_q}) as follows:
\end{ntasks}
\end{region}

\begin{region}[x=74.7,right=0]{379.3}
\begin{lstlisting}[style=py]
elif event.type == pygame.KEYDOWN:
    if event.key == pygame.K_q:
        running = False
    elif event.key == pygame.K_UP:
        rockets.append(Rocket(Paddle.x + Paddle.w//2, Paddle.y, -10))
\end{lstlisting}
\end{region}

\begin{region}[x=74.7,right=0]{451.4}
The start position is the middle of the paddle, and the velocity is $-10$
(because our y-axis points downward).
\end{region}

\begin{region}{488.3}
\begin{ntasks}[start=4]
\item \textbf{Moving and drawing the rockets}\\
  So far, the main loop moves and draws all balls -- now you do the same for
  the rockets. Directly afterwards:
\end{ntasks}
\end{region}

\begin{region}[x=74.7,right=0]{536.6}
\begin{lstlisting}[style=py]
for r in rockets:
    r.draw()
    r.move()
\end{lstlisting}
\end{region}

\begin{region}[x=74.7,right=0]{584.9}
From now on, you can test the feature with \code{python breakout.py}.
\end{region}

\begin{region}[pitch=12.3]{609.9}
\begin{ntasks}[start=5]
\item \textbf{Once all this runs and the rockets fly -- into the repository:}\\
  \code{git add breakout.py rocket.py}\\
  \code{git commit -m "added rockets"}\\
  \code{git push}
\end{ntasks}
\end{region}

\nextpage
%% ==================================================================== 14 ==
\begin{region}{69.4}
\heading{Step 9: Finally, we clean up!}
\end{region}

\begin{region}[pitch=14]{97.0}
Our rockets stick to the ceiling -- and on purpose!\\
So far, the method \code{move} only checked whether \code{self.y < 10} and
then simply \textbf{did nothing more}. This is why they stay right there.\\
The value \code{10} was \textbf{deliberately} chosen instead of \code{0} --
this way, you can see the ``garbage'' that our code has produced so far.
\end{region}

\begin{region}[pitch=13.53]{180.3}
\subheading{Task}
Now we clean up: the class \code{Rocket} gets a new attribute \code{alive}.
When a rocket reaches the ceiling (i.e., \code{y < 10}), this attribute is set
to \code{False} -- and we then remove the rocket from the list
\code{rockets}.
\end{region}

\begin{region}{248.3}
\textbf{New class Rocket in \codemd{rocket.py}:}
\end{region}

\begin{region}{261.1}
\begin{lstlisting}[style=pysmall]
class |\cy{Rocket}|:
\end{lstlisting}
\end{region}

\begin{region}[pitch=11]{273.2}
\begin{lstlisting}[style=pysmall]
    def |\cy{\_\_init\_\_}|(self, x, y, vy):
        self.x = x
        self.y = y
        self.vy = vy
        self.c = (255, 255, 255)
        self.alive = True

    def |\cy{move}|(self):
        if self.y < 10:
            self.alive = False
            return
        self.y += self.vy

    def |\cy{draw}|(self):
        pygame.draw.circle(Game.screen, self.c, (self.x, self.y), 5)
\end{lstlisting}
\end{region}

\begin{region}[pitch=13.3]{453.7}
\textbf{In the main loop of \codemd{breakout.py}}\\
add the following line after drawing and moving the rockets:
\end{region}

\begin{region}{493.6}
\begin{lstlisting}[style=py]
    rockets = [item for item in rockets if item.alive]
\end{lstlisting}
\end{region}

\begin{region}[pitch=14.3]{519.4}
\subheading{What does this line do?}
This is a so-called \textbf{list comprehension}. It creates a new list that
\textbf{contains only the rockets} for which \code{alive == True} -- that is,
those that have not yet ``exploded'' at the ceiling. \textbf{In short:}
``Keep only the living rockets.''
\end{region}

\begin{region}[pitch=13.6]{589.7}
\textbf{If everything works}, you can improve the \code{move} method of a
rocket a little: set \code{alive} to False as soon as the y-coordinate is less
than \code{0}.
\end{region}

\begin{region}[pitch=13.4]{630.6}
\textbf{Attention:} Afterwards, be sure to test once more whether everything
still works correctly -- to rule out typical last-minute typos.
\end{region}

\begin{region}[pitch=13.35]{671.2}
\subheading{Then into the repository:}
\code{git add rocket.py breakout.py}\\
\code{git commit -m "rockets disappear at the ceiling"}\\
\code{git push}
\end{region}

\begin{region}[pitch=13.2]{738.6}
\subheading{Final remark:}
In the same way, you can also make balls disappear that you miss with the
paddle. That is where we will start next time -- to pick up right where we
left off today.
\end{region}

\end{document}
